\section{Peta Pencapaian 19 Sub-CPMK}
\label{sec:ch17-peta}

\subsection{Rekapitulasi}

Seluruh sembilan belas Sub-CPMK telah dibahas pada bab-bab inti.
Tabel berikut merangkum bab utama, asesmen, dan bukti belajar yang
diharapkan.

{\small
\begin{longtable}{@{}p{1.2cm}p{7.0cm}p{1.6cm}p{4.4cm}@{}}
\caption{Peta pencapaian Sub-CPMK}
\label{tab:ch17-peta}\\
\toprule
\textbf{Kode} & \textbf{Kompetensi} & \textbf{Bab} &
\textbf{Bukti belajar} \\
\midrule
\endfirsthead
\toprule
\textbf{Kode} & \textbf{Kompetensi} & \textbf{Bab} &
\textbf{Bukti belajar} \\
\midrule
\endhead
\bottomrule
\endfoot
1.1 & Membedakan interpreter, compiler, dan translator & III &
Kuis, ujian tulis \\
1.2 & Menjelaskan fase kompilasi, tabel simbol, dan error handling &
III & Kuis, diagram pipeline \\
1.3 & Menganalisis one-pass/multi-pass dan run-time environment &
III & Tugas teori, analisis kasus \\
1.4 & Menjelaskan hierarki Chomsky dan automata & IV &
Tugas klasifikasi, UTS \\
1.5 & Menjelaskan attribute grammar, SDD, dan translation scheme &
XIV & Tugas teori, P4 \\
1.6 & Mengimplementasikan symbol table dan type checking &
XV & P4, kode dan tes \\
1.7 & Menangani semantic error dengan pesan informatif &
XVI & P4, proyek akhir \\
2.1 & Menulis ekspresi reguler untuk token & V &
Tugas teori, P1-A \\
2.2 & Mengimplementasikan NFA/DFA untuk token recognition &
VI & P1-B, kode \\
2.3 & Menjelaskan ekivalensi NFA-DFA dan minimalisasi & VI &
Tugas teori, UTS \\
2.4 & Mengimplementasikan lexical analyzer & VII &
P1, kode dan tes \\
3.1 & Menulis CFG, BNF/EBNF, derivasi, dan parse tree & VIII &
P2-A, ujian tulis \\
3.2 & Menganalisis ambiguitas, left recursion, left factoring &
IX & Tugas teori, UTS \\
3.3 & Menghitung FIRST/FOLLOW dan menjelaskan LL(1)/LR(1) &
X & UTS, P2 \\
3.4 & Merancang CFG untuk bahasa sederhana & VIII &
P2-A, proyek \\
4.1 & Membangun recursive descent dan predictive parser & X &
P2, kode dan tes \\
4.2 & Menerapkan bottom-up parsing dan menganalisis konflik &
XI--XII & UTS, P3 \\
4.3 & Menggunakan parser generator dan semantic actions & XIII &
P3, kode \\
4.4 & Membangun AST dan memvalidasi kode & XIII, XVI &
P3, proyek akhir \\
\end{longtable}
}

\subsection{Analisis Ketercapaian}

Ketercapaian dapat dianalisis pada tiga tingkat:
\begin{enumerate}[label=\arabic*.,leftmargin=*]
  \item \textbf{pengetahuan}: mahasiswa dapat menjelaskan konsep;
  \item \textbf{aplikasi}: mahasiswa dapat menerapkan teknik;
  \item \textbf{integrasi}: mahasiswa dapat menggabungkan teknik
        dalam proyek.
\end{enumerate}

Nilai kuis dan ujian mengukur tingkat pengetahuan dan aplikasi.
Praktikum dan proyek mengukur aplikasi dan integrasi. Kombinasi
keduanya memberikan gambaran kompetensi yang lebih lengkap.

\begin{figure}[H]
\centering
\begin{tikzpicture}[node distance=6mm]
  \node[phase, text width=3.4cm] (know) {Pengetahuan\\kuis, ujian};
  \node[phase, text width=3.4cm, right=of know] (apply)
        {Aplikasi\\praktikum};
  \node[phase, text width=3.4cm, right=of apply] (integrate)
        {Integrasi\\proyek akhir};
  \draw[arrow] (know) -- (apply);
  \draw[arrow] (apply) -- (integrate);
\end{tikzpicture}
\caption{Tiga tingkat ketercapaian kompetensi}
\label{fig:ch17-levels}
\end{figure}
